\documentclass{article}
\usepackage{amsmath,amssymb}
\usepackage{xurl}
\usepackage[hidelinks]{hyperref}
% Shared commands for notes in docs/paper-gaps/.

\DeclareUnicodeCharacter{2080}{\ifmmode{}_{0}\else\textsubscript{0}\fi}
\DeclareUnicodeCharacter{2081}{\ifmmode{}_{1}\else\textsubscript{1}\fi}
\DeclareUnicodeCharacter{2082}{\ifmmode{}_{2}\else\textsubscript{2}\fi}
\DeclareUnicodeCharacter{2099}{\ifmmode{}_{n}\else\textsubscript{n}\fi}

\newcommand{\Lean}{\textsc{Lean}}
\ExplSyntaxOn
\NewDocumentCommand{\leanid}{m}{
  \ifmmode
    \textsf{\detokenize{#1}}
  \else
    \group_begin:
    \urlstyle{sf}
    \tl_set:Nn \l_tmpa_tl {#1}
    \tl_replace_all:Nnn \l_tmpa_tl {\_} {_}
    \exp_args:NV \path \l_tmpa_tl
    \group_end:
  \fi
}
\ExplSyntaxOff
\providecommand{\tr}{\operatorname{tr}}
\newcommand{\tnleanIssueBase}{https://github.com/LionSR/TNLean/issues/}
\newcommand{\tnleanPrBase}{https://github.com/LionSR/TNLean/pull/}
\newcommand{\tnleanDocBase}{https://sirui-lu.com/TNLean/docs/}
\newcommand{\ghissue}[1]{\href{\tnleanIssueBase#1}{GitHub issue \##1}}
\newcommand{\ghpr}[1]{\href{\tnleanPrBase#1}{PR \##1}}
\newcommand{\doclink}[1]{\href{\tnleanDocBase#1.html}{\path{#1}}}

% Dirac notation and the identity operator, as in blueprint/src/macros/common.tex.
% \providecommand keeps note-local definitions authoritative.
\usepackage{dsfont}
\providecommand{\ket}[1]{|#1\rangle}
\providecommand{\bra}[1]{\langle #1|}
\providecommand{\braket}[2]{\langle #1 | #2 \rangle}
\providecommand{\ketbra}[2]{|#1\rangle\!\langle #2|}
\providecommand{\Id}{\mathds{1}}
\providecommand{\one}{\mathds{1}}
\providecommand{\C}{\mathbb{C}}
\providecommand{\MN}[1]{M_{#1}(\C)}
\providecommand{\MD}{M_D(\C)}

% Verdict marker: \gapnote{<kind>}{<status>}. Kind is the policy
% classification (clarification, local-correction, scope-restriction,
% unfaithful, false-source, open-gap); status is open, wip (elimination
% underway), resolved, or historical. Machine-read by the site index;
% prints a verdict line.
\newcommand{\gapnote}[2]{\par\noindent\textbf{Verdict:} #1 (#2).\par}

\title{Commuting Parent Hamiltonians of $G$-Isometric Tensors: the Projector and the Scope}
\author{}
\date{2026-10-05}
\begin{document}
\maketitle
\gapnote{scope-restriction}{open}

\section{What the source prints}
Section~6.4 of \cite{Schuch2010PEPS} shows that the parent Hamiltonians of
$G$-isometric PEPS are sums of commuting local terms, and announces that the
proof is given for matrix product states, the generalization to PEPS being
straightforward.\footnote{Source traceability:
    \path{Papers/1001.3807/paper_v3.tex}, lines 2092--2096.}
For a $G$-isometric matrix product tensor $A$, on whose bond the
left-regular representation acts, let
\begin{align}
    \mathcal S_2=\Bigl\{\sum_{ij}\operatorname{tr}[A^iA^jX]\,\ket{ij}
        \Bigm| X\Bigr\}.
    \label{eq:scp10_s2}
\end{align}
The lemma of that section states that the local terms $h_i$ of the parent
Hamiltonian, ``which project onto the complement of $\mathcal S_2$'', are of
the form of equation ``eq:iso:ham-proj-from-A'', the two-site operator
\begin{align}
    P=\sum_{ijkl}\operatorname{tr}\bigl[A^iA^j(A^kA^l)^\dagger\bigr]
        \ket{ij}\bra{kl}
    \label{eq:scp10_projector}
\end{align}
drawn in figure ``ham-proj-from-A''. Its proof shows that $P$ is a projector,
that its range lies in $\mathcal S_2$, that it fixes every vector of
$\mathcal S_2$, and that it is self-adjoint (figures
``ham-proj-from-A-is-proj'' and ``ham-proj-from-A-pres-subspace'').
The theorem that follows states that for $G$-isometric PEPS the terms $h_i$
of the parent Hamiltonian of the two-dimensional parent-Hamiltonian theorem
commute. Its proof writes the products of the operators $P$ on sites $1,2$
and $2,3$ as the two sides of figure ``ham-comm-step1'', which are equal
because the sum over $g$ of the virtual operators $U_g\otimes U_g^{-1}$ can
be moved from one bond to the other, and closes with ``the same proof applies
to two dimensions''.\footnote{Source traceability:
    \path{Papers/1001.3807/paper_v3.tex}, lines 2098--2153, figures
    \path{Papers/1001.3807/figs4/ham-proj-from-A.pdf},
    \path{Papers/1001.3807/figs4/ham-proj-from-A-is-proj.pdf},
    \path{Papers/1001.3807/figs4/ham-proj-from-A-pres-subspace.pdf},
    \path{Papers/1001.3807/figs4/ham-comm-step1.pdf} and
    \path{Papers/1001.3807/figs4/ham-comm-step2.pdf}.}

\section{Local fix: the operator is the projector onto $\mathcal S_2$}
The local term of the parent Hamiltonian is
$h=1-\Pi_{\mathcal S}$.\footnote{Source traceability:
    \path{Papers/1001.3807/paper_v3.tex}, lines 1177--1183 and 1534--1539.}
The proof of the lemma shows that $P$ is the orthogonal projector onto
$\mathcal S_2$, not onto its complement, so the operator of
equation ``eq:iso:ham-proj-from-A'' is $1-h_i$. The commutation theorem is
unaffected: $h_1$ and $h_2$ commute exactly when $1-h_1$ and $1-h_2$ do.
In addition the source omits normalizations in diagrams. For a tensor $A$
with matrices $A^i$ on the bond space $\mathbb C[G]$, let
$\mathcal P(A)X=\sum_i\operatorname{tr}[A^iX]\ket i$ be the map from the
virtual operators $X$ to the physical system of one site, let
$\mathcal P(A^\dagger)\ket i=(A^i)^\dagger$ be the corresponding map for the
tensor $A^\dagger$, and let
$\sigma(X)=|G|^{-1}\sum_{g\in G}L_gXL_g^{-1}$ be the twirl onto the
commutant of the left-regular representation. The formal notion of
$G$-isometry allows $\mathcal P(A^\dagger)\mathcal P(A)=c\,\sigma$
with one constant $c>0$ \cite{gap:rmp_peps_quantum_double_g_isometry}; then
$P=c_2\,\Pi_{\mathcal S_2}$ for the constant $c_2>0$ of the two-site tensor
$A^iA^j$. The formalization therefore states
\begin{align}
    h=1-c_2^{-1}P,
    \label{eq:scp10_local_term}
\end{align}
together with $(c_2^{-1}P)^2=c_2^{-1}P$, $\operatorname{range}P=\mathcal S_2$
and $P^\dagger=P$. The last identity holds for every tensor, since the
matrix entries of $P$ satisfy
$\overline{\operatorname{tr}[A^kA^l(A^iA^j)^\dagger]}
    =\operatorname{tr}[A^iA^j(A^kA^l)^\dagger]$. The source's closing remark,
that with $A^{-1}$ in place of $A^\dagger$ the operator is a projector onto
$\mathcal S_2$ that is not self-adjoint, is not formalized.

\section{The proof of the commutation}
Applying $\mathcal P(A^\dagger)\mathcal P(A)=c\,\sigma$ to an operator $N$
that commutes with every $L_g$, for which $\sigma(N)=N$, gives
\begin{align}
    \sum_b\operatorname{tr}[A^bN]\,(A^b)^\dagger=c\,N .
    \label{eq:scp10_isometry_identity}
\end{align}
The tensor $A$ is invariant, $L_gA^iL_g^{-1}=A^i$, and since $L_g$ is unitary
in the basis of group elements, $L_g^\dagger=L_{g^{-1}}$, each adjoint
$(A^i)^\dagger$ also commutes with every $L_g$. Every product of the
matrices $A^i$ and $(A^i)^\dagger$ therefore commutes with every $L_g$; this
is the step in which the sum over $g$ of figure ``ham-comm-step1'' drops out.
Multiplying \eqref{eq:scp10_isometry_identity} on the left by any operator
$M$ and taking the trace gives the trace pairing
\begin{align}
    \sum_b\operatorname{tr}[M(A^b)^\dagger]\,\operatorname{tr}[A^bN]
    =c\operatorname{tr}[MN].
    \label{eq:scp10_trace_pairing}
\end{align}
The matrix entry of $(P\otimes1)(1\otimes P)$ at the configurations
$s_0s_1s_2$ and $w_0w_1w_2$ is computed with
$M=(A^{w_0})^\dagger A^{s_0}A^{s_1}$ and
$N=A^{s_2}(A^{w_2})^\dagger(A^{w_1})^\dagger$, which commutes with every
$L_g$. By cyclicity of the trace, and then by
\eqref{eq:scp10_trace_pairing},
\begin{align}
    \sum_b\operatorname{tr}\bigl[A^{s_0}A^{s_1}(A^{w_0}A^b)^\dagger\bigr]\,
        \operatorname{tr}\bigl[A^bA^{s_2}(A^{w_1}A^{w_2})^\dagger\bigr]
    &=\sum_b\operatorname{tr}[M(A^b)^\dagger]\,\operatorname{tr}[A^bN]
    \nonumber\\
    &=c\operatorname{tr}[MN]
    =c\,\operatorname{tr}\bigl[A^{s_0}A^{s_1}A^{s_2}
        (A^{w_0}A^{w_1}A^{w_2})^\dagger\bigr].
    \label{eq:scp10_three_site_entry}
\end{align}
For $(1\otimes P)(P\otimes1)$ the same computation applies with
$M=A^{s_1}A^{s_2}(A^{w_2})^\dagger$ and
$N=(A^{w_1})^\dagger(A^{w_0})^\dagger A^{s_0}$.
The two matrices agree, so both products are
$c$ times the three-site operator of the same form. The two translated terms
on any three consecutive sites of a ring commute, and terms with disjoint
supports commute trivially.

\section{Two-dimensional regional projectors}
For a finite simple graph, assign a group label $\alpha_{v,e}$ to each
half-edge. The actual canonical regular open-region range consists exactly
of functions invariant under independent left multiplication at the region
vertices and simultaneous right multiplication at the ends of each internal
edge, and vanishing unless there is a vertex assignment $k$ with
$k_v^{-1}\alpha_{v,e}=k_w^{-1}\alpha_{w,e}$ on every internal edge $e=(v,w)$.
This equality is proved from the original open-region contraction, including
all independent dangling boundary labels; neither a rank formula nor
connectedness is assumed.

Let $A_v$ average the left action at vertex $v$, $B_e$ average the shared
right action on edge $e$, and $D_R$ be the diagonal indicator of this flat
locus on a region $R$. The extended actual regional range projector is
\[
  Q_R=\left(\prod_{v\in R}A_v\right)
       \left(\prod_{e\subseteq R}B_e\right)D_R.
\]
Left and right multiplication commute. Averages on distinct vertices or
edges act on distinct half-edge coordinates; coincident averages are
identical. Flatness is invariant under every left vertex action and shared
right edge action, and the $D_R$ are diagonal. Thus $Q_RQ_S=Q_SQ_R$ for
arbitrary overlapping regions.

For the original physical site maps, local $G$-isometry gives
$A^\dagger A=cE$ with $c>0$ and $E$ the regular invariant projector.
A supported orthogonal projector $Q$ is transported to
$c^{-1}AQA^\dagger$, whose range is exactly the original physical regional
range. Its intertwining relation is $PA=AQ$.
No surjectivity onto the ambient physical space is imposed. The normalized
local adjoints give nonzero physical site-image projectors $J_v$; nonzero
follows from the constant vector in the regular invariant subspace.

The physical regional projectors commute with every product of site-image
projectors. For regions $R,S$, their commutator is supported on $R\cup S$
and absorbs the product of site-image projectors there. Its product with
the full product of site maps is zero by the virtual commutation. Multiplying
by the product of normalized adjoints and canceling the nonzero complementary
site-image factor therefore makes the commutator zero on the entire physical
space, including unused local physical directions. The parent interactions
are the complements of these range projectors and consequently commute.\footnote{
The source-facing statement is
\texttt{regularIsometric\_canonicalRegionParentInteractions\_commute} in
\path{TNLean/PEPS/ParentHamiltonian/GIsometricRegionParentHamiltonian.lean}.
Its hypotheses are local regular $G$-isometry; all support, nonzero-image,
and regional intertwining identities are derived.}

\section{Scope and verdict}
The two-dimensional commutation assertion is now proved for the original
canonical physical parent interactions on every finite simple graph and
arbitrary overlapping regions. In particular, it covers plaquette terms on
ordinary square tori represented by this graph convention. The source's
projector-complement error is corrected as in
\eqref{eq:scp10_local_term}.

A residual indexing restriction remains: the simple-graph description does
not encode self-loops or parallel bonds in degenerate periodic presentations.
The older matrix-product theorem still states $N\geq3$ and does not by itself
cover the two-site periodic double-bond convention. Square tori with both
periods at least three have no such issue. No assertion about those smaller
periodic presentations is inferred from the finite-simple-graph theorem.
Eliminating this last restriction requires an incidence description retaining
multiple half-edge occurrences, or a direct proof for the exceptional
periodic contractions; it does not require another two-dimensional range or
physical-transport hypothesis.

\bibliographystyle{alpha}
\bibliography{references}
\end{document}
